% Falsilla (hoja guía para poner debajo) (A4) — Eric Lucero González — https://ericlucero.dev
% Licencia CC BY-NC-SA 4.0. Compilar con pdfLaTeX (dos pasadas: latexmk lo hace solo).
\documentclass[a4paper]{article}
\usepackage[T1]{fontenc}
\usepackage[spanish]{babel}
\usepackage{lmodern}
\usepackage{tikz}
\pagestyle{empty}

% Líneas gruesas negras pensadas para verse por transparencia debajo de una hoja blanca.
% ── Configuración (mm) ─────────────────────────────────────
\def\anchoPapel{210}
\def\altoPapel{297}
\def\margenLateral{25}   % posición de las líneas de margen, desde cada borde
\def\margenVertical{25}  % margen mínimo arriba y abajo
\def\paso{9}            % distancia entre líneas
\def\sobresale{10}      % cuánto pasan las líneas de margen del primer y último renglón

\begin{document}
\begin{tikzpicture}[
    remember picture, overlay, shift={(current page.south west)},
    x=1mm, y=1mm,
    master lines/.style={black, ultra thick},
    margin lines/.style={black!70, thick, dashed}
  ]
  % Renglones que caben, centrados en vertical
  \pgfmathtruncatemacro{\n}{(\altoPapel - 2*\margenVertical) / \paso}
  \pgfmathsetmacro{\yIni}{(\altoPapel - \n*\paso) / 2}
  \pgfmathsetmacro{\yFin}{\yIni + \n*\paso}
  \pgfmathsetmacro{\xIzq}{\margenLateral}
  \pgfmathsetmacro{\xDer}{\anchoPapel - \margenLateral}

  \foreach \i in {0,...,\n} {
    \draw[master lines] (\xIzq, \yIni + \i*\paso) -- (\xDer, \yIni + \i*\paso);
  }

  \draw[margin lines] (\xIzq, \yIni - \sobresale) -- (\xIzq, \yFin + \sobresale);
  \draw[margin lines] (\xDer, \yIni - \sobresale) -- (\xDer, \yFin + \sobresale);
\end{tikzpicture}
\end{document}
